\documentclass{tufte-handout} % A4 paper and 11pt font size
\usepackage[activate={true,nocompatibility},final,tracking=true,kerning=true,spacing=true,factor=1100,stretch=10,shrink=10]{microtype}
\usepackage[T1]{fontenc} % Use 8-bit encoding that has 256 glyphs
\usepackage{mathpazo} % Use the Adobe Utopia font for the document - comment this line to return to the LaTeX default
\usepackage[english]{babel} % English language/hyphenation
\usepackage{amsmath,amsfonts,amsthm, amssymb} % Math packages
\usepackage{pgf,tikz}
\usetikzlibrary{positioning,matrix,arrows}
\usepackage{float}
\usepackage{tikz-cd}
\usepackage{caption}
\usepackage{stmaryrd}
\usepackage{multicol}
\usepackage{booktabs}
\usepackage{verbatim}
\usepackage{lipsum} % Used for inserting dummy 'Lorem ipsum' text into the template
\usepackage{sectsty} % Allows customizing section commands
\allsectionsfont{\normalfont \bfseries} % Make all sections centered, the default font and small caps
\usepackage{enumerate}
\usepackage{pythonhighlight}
\usepackage{fancyhdr} % Custom headers and footers
\pagestyle{fancyplain} % Makes all pages in the document conform to the custom headers and footers
\fancyhead{} % No page header - if you want one, create it in the same way as the footers below
\fancyfoot[L]{} % Empty left footer
\fancyfoot[C]{} % Empty center footer
\fancyfoot[R]{\thepage} % Page numbering for right footer
\renewcommand{\headrulewidth}{0pt} % Remove header underlines
\renewcommand{\footrulewidth}{0pt} % Remove footer underlines
\setlength{\headheight}{13.6pt} % Customize the height of the header
\allowdisplaybreaks
%--------Theorem Environments--------
%theoremstyle{plain} --- default


\def \v {\vspace{0.2cm}}

\geometry{
	left=13mm, % left margin
	textwidth=130mm, % main text block
	marginparsep=8mm, % gutter between main text block and margin notes
	marginparwidth=55mm % width of margin notes
}
\fontsize{10}{20}\selectfont


%\numberwithin{equation}{section}

\setcounter{secnumdepth}{2}

\theoremstyle{theorem}
\newtheorem{thm}{Theorem}[subsection]
\newtheorem{lem}[thm]{Lemma}

\theoremstyle{definition}
\newtheorem{defn}[thm]{Definition}

\theoremstyle{remark}
\newtheorem{rem}[thm]{Remark}

\theoremstyle{remark}
\newtheorem{exa}[thm]{Example}
%----------------------------------------------------------------------------------------
%	TITLE SECTION
%----------------------------------------------------------------------------------------

\title{	
	\normalfont\normalsize 
	{University of Cape Town} \\ [0pt] % Your university, school and/or department name(s)
	\huge List of Theorems and Definitions% The assignment title
}\author{KB} % Your name
\date{\vspace{-5pt}\normalsize\today} % Today's date or a custom date

\begin{document}
	\justifying 
	\maketitle

This is just a summary of the definitions and theorems used.
The theorems are paraphrased, and not the actual theorems stated.
This is a growing list.
	
\section{Monads}
\subsection{Monads and their Algebras}


		Note on composeability of functors and natural transformations
		
		\begin{defn} Monad
		\end{defn}
		\begin{defn} $T$-algebra
		\end{defn} $\eta$ and $\mu$ of a monad $T$ are referred to as the identity and multiplication of a monad
		\begin{defn} Morphism of $T$-algebras
		\end{defn}
		\begin{defn} Eilenberg-Moore category of $ T $, written $\mathbb{X}^T $
		\end{defn}
		\begin{thm} Every adjunction determines a monad. \textit{Proof included.}
		\end{thm}
		\begin{defn} Monad defined by the adjunction
		\end{defn}
		\begin{thm} Every monad is defined by its $ T $-algebras: every monad determines an adjunction and the monad defined by this adjunction is the same as the first monad. \textit{Proof included, with full description of the adjunction.}
		\end{thm}
		\begin{thm} Comparison of adjunctions with algebras: existence of Kleisli comparison functor. \textit{Proof included, full description of functor.}
		\end{thm}
		\begin{thm} Existence of adjunction between first category $ \mathbb{A} $ and $\mathbb{X}^T $. \textit{Proof included, with full definition of adjunction.}
		\end{thm}
		\begin{defn} Comparison functor
		\end{defn}
		\begin{defn} Monadic functor
		\end{defn}
		\begin{thm} \textcolor{blue}{Kleisli category} of a monad. \textit{Proof included.}
		\end{thm}
		\begin{thm} \textcolor{blue}{Comparison theorem} for the Kleisli construction.
		\end{thm}
		\begin{thm} The \textcolor{blue}{Eilenberg-Moore and Kleisli} construction as terminal and initial objects.
		\end{thm}
		
\subsection{Beck's Theorem}

		\begin{defn}
			\textcolor{blue}{Coequaliser, fork, split coequaliser, split fork, absolute coequaliser, reflexive coequaliser}
		\end{defn}	
		\begin{thm}
			\textcolor{blue}{Beck's theorem characterising algebras:} the 3 equivalent conditions about $ K $ being an equivalence, and how the right adjoint creates coequalisers. \textit{Proof included.}
		\end{thm}
		\begin{defn}
			A functor reflects isos
		\end{defn}
		\begin{thm}
			Beck's monadicity theorem: 4 equivalent conditions of monadicity, in terms of preserving coeqaliser diagrams of a given form. \textit{Proof included.}
		\end{thm}
	    \begin{thm}
			Crude monadicity theorem: in which a functor is monadic if it has a left adjoint, reflects isos and preserves reflexive coequalisers. \textit{Proof included}
		\end{thm}

\subsection{Comonads and Precomonadicity}
	    \begin{defn}
			\textcolor{blue}{Comonad}
		\end{defn}
		\begin{defn}
			\textcolor{blue}{Precomonadic}
		\end{defn}

\end{document}